\section{Conclusion}

The research presented in this article introduces the CLIP-Guided Dual-Path Diffusion Model (CDDM), a pioneering cross-modal generative framework designed to address the formidable challenges of zero-shot fault diagnosis (ZSFD) in complex industrial environments. By bridging the gap between high-level textual diagnostic descriptions and low-level industrial time-series signals, CDDM provides a robust solution for detecting unseen fault categories without the need for historical labeled samples. The integration of large language models (LLMs) and contrastive language-image pre-training (CLIP) architectures enables the mapping of expert-level fault knowledge into a continuous semantic space, effectively overcoming the limitations of traditional discrete binary attribute matrices.

\subsection{Summary of Contributions}
The technical contributions of this work are fourfold, each addressing a critical bottleneck in existing ZSFD methodologies.

Firstly, the development of LLM-based continuous semantic embeddings represents a significant departure from conventional attribute-based learning. Traditional models rely on hard-coded, sparse binary matrices that often fail to capture the nuanced dependencies between fault causes and their functional effects. By leveraging the semantic richness of LLMs, we transform discrete expert descriptions into a dense, continuous manifold. This innovation not only preserves fine-grained semantic information but also enhances the model's ability to generalize across diverse fault domains, virtually eliminating the reliance on manually annotated attribute tables that are prone to human error and information loss.

Secondly, the CDDM introduces a generative dual-path diffusion architecture tailored for industrial signal synthesis. Unlike GAN-based approaches that suffer from mode collapse and training instability, our dual-path design—incorporating a spatial-temporal path (1D-CNN) and a time-frequency spectral path (attention-based)—ensures the generation of high-fidelity, physically plausible signals. This structural bifurcation allows the model to simultaneously capture local transient features and global periodic characteristics inherent in industrial sensor data, providing a specialized backbone for robust zero-shot knowledge transfer.

Thirdly, the integration of a Multi-Modal Contrastive Consistency Loss serves as a pivotal alignment mechanism. By enforcing semantic-structural agreement between the generated signal representations and the guiding textual embeddings, this loss function ensures that the synthesized unseen fault samples are not merely realistic but are strictly conditioned on the intended fault physics. This replaces standard attribute regression techniques and provides an explicit mathematical guarantee against the domain shift problem, which has historically plagued generative ZSL models.

Finally, extensive experimental validation across three benchmark datasets—the Three-Phase Transmission System (TPTS), the Tennessee Eastman Process (TEP), and a complex hydraulic system—underscores the practical efficacy of CDDM. In the TEP case study, CDDM achieved a state-of-the-art average diagnostic accuracy of 85.86\%, outperforming GAN-based (FAGAN: 73.89\%) and VAE-based (VAEGAN-AR: 81.61\%) baselines. Similarly, in the hydraulic system benchmark, CDDM maintained superior performance even as the number of unseen classes increased to 30, recording a significant margin of 8.88\% over the nearest diffusion-based competitor. These results confirm that CDDM provides a more reliable and scalable foundation for generalized condition monitoring.

\subsection{Limitations and Future Work}
Despite the substantial advancements demonstrated by the CDDM framework, several limitations offer fertile ground for future investigation. A primary constraint lies in the iterative nature of the denoising diffusion process. While CDDM produces superior sample fidelity, the multi-step reverse sampling phase introduces a computational bottleneck that may impede its deployment in strictly real-time diagnostic scenarios on resource-constrained edge devices. Industrial applications requiring microsecond-level response times for fault isolation may find the current inference latency suboptimal. Consequently, future research will explore accelerated sampling strategies, such as Denoising Diffusion Implicit Models (DDIM \cite{song2021denoising}), consistency models, and knowledge distillation techniques to compress the reverse chain without sacrificing generative quality.

Another limitation concerns the model's dependence on the quality and objectivity of the initial textual fault descriptions. While LLMs mitigate discrete attribute errors, the "expert knowledge" provided as input remains a potential source of bias. If the textual descriptions are overly generic or fail to mention critical transient behaviors, the generated signals may lack the discriminative features necessary for accurate classification. To address this, we aim to develop an "Active Semantic Discovery" module that can autonomously refine or suggest expansions to expert descriptions by cross-referencing historical seen-class data and industrial ontologies.

Furthermore, while our dual-path architecture effectively handles 1D sensor signals and 2D spectrograms, real-world machinery often relies on a broader array of heterogeneous data, including transient acoustic emissions, video feeds of mechanical wear, and full-field thermal imaging. Extending the current framework to accommodate true multi-modal sensor fusion beyond simple concatenation is a critical next step. We envisage a "Hyper-Graph Diffusion" approach where different sensor modalities are treated as nodes in a conditioned graph, allowing the model to learn the complex inter-dependencies between, for example, a pressure drop and a corresponding thermal spike.

Finally, we recognize a performance bottleneck in certain extreme "out-of-distribution" scenarios. When an unseen fault deviates radically from the physical principles observed in the seen classes, the semantic-to-signal mapping may suffer from extrapolation inaccuracies. Future work will investigate the integration of physics-informed neural networks (PINNs) into the diffusion backbone to ensure that even the most "novel" generated signals remain bounded by known laws of thermodynamics and fluid mechanics.

\subsection{Industrial Implications}
The deployment of CDDM carries profound implications for the advancement of Industry 4.0 and the realization of "lights-out" manufacturing. In safety-critical sectors such as nuclear power, aerospace, and large-scale chemical processing, certain catastrophic faults are too dangerous or costly to reproduce experimentally. CDDM's ability to synthesize high-fidelity training data for these rare events directly from textual specifications allows for the preemptive training of deep diagnostic classifiers. This shifts the industrial maintenance paradigm from "reactive/preventative" to "proactively intelligent," drastically reducing unplanned downtime and the associated multi-million dollar costs.

Moreover, CDDM offers immense value during the commissioning of new industrial facilities. Traditionally, diagnostic systems require months of data collection to reach reliable accuracy levels. By utilizing the CDDM framework, engineers can "bootstrap" the diagnostic system on Day 1 using digital twin specifications and expert manuals, providing an immediate layer of protection for new equipment. The potential for integration with Industrial Internet of Things (IIoT) platforms and Digital Twin ecosystems is significant; CDDM could serve as the "generative engine" that populates a digital twin's failure mode library with synthetic yet physically grounded examples, enabling more realistic stress-testing of control logic and maintenance protocols.

\subsection{Concluding Remarks}
In conclusion, the CLIP-Guided Dual-Path Diffusion Model (CDDM) represents a transformative step in industrial zero-shot fault diagnosis. By synergizing the interpretability of textual expert knowledge with the generative power of diffusion models, we have established a framework that not only surpasses the accuracy of current state-of-the-art models but also provides a more robust and scalable solution for real-world deployment. As industrial systems continue to grow in complexity and data volume, the ability to "diagnose the unseen" through cross-modal intelligence will be the cornerstone of resilient and autonomous machinery operation. This work lays the foundation for a new generation of diagnostic systems that can reason, synthesize, and protect modern industrial infrastructure with unprecedented foresight.

\cite{lee2014prognostics, wang2020deep, precup2015overview, tzeng2017adversarial, song2021denoising}
